\documentclass[11pt,a4paper]{article}
\usepackage[margin=2.3cm]{geometry}
\usepackage{amsmath,amssymb,booktabs,graphicx,microtype}
\usepackage[colorlinks=true,linkcolor=blue!50!black,citecolor=blue!50!black,urlcolor=blue!50!black]{hyperref}
\newcommand{\eps}{\epsilon}
\newcommand{\dlog}{\mathrm{d}\!\log}
\title{Direct numerical solution of the canonical differential equation\\ for the family $H$ (Euclidean and continued region)}
\author{implementation notes}
\date{\today}
\begin{document}
\maketitle

\section{Problem}
The $N_f=371$ canonical master integrals $\vec J(x,y;\eps)$ of the family $H$ satisfy
\begin{equation}
  \mathrm d\vec J=\eps\,A\,\vec J,\qquad A=\sum_{k\in\mathcal A} a_k\,\dlog l_k(x,y),\qquad z\equiv 1-x-y ,
\end{equation}
with constant $371\times371$ matrices $a_k$ (\texttt{atilde.m}) and the 14 active letters
$\mathcal A=\{1,\dots,13,16\}$:
\begin{gather*}
 l_1=x,\ l_2=y,\ l_3=z,\ l_4=1-x,\ l_5=1-y,\ l_6=x+y,\ l_9=1-(1+x)y,\ l_{10}=1-x(1+y),\\
 l_{11}=x-y+y^2,\ l_{12}=-x+x^2+y,\ l_{13}=(1-x)^2-y,\ l_{16}=(1-y)^2-x,\\
 l_7=\frac{xz-i r}{xz+i r},\qquad l_8=\frac{xy-i r}{xy+i r},\qquad r^2=b\equiv xyz .
\end{gather*}
The two odd letters contain the only square root.  Using $r^2=b$ (so $r\,\mathrm dr=\mathrm db/2$) their
$\dlog$ is written with a single overall $1/r$:
\begin{equation}\label{eq:odd}
 \dlog\frac{a-ir}{a+ir}=\frac{i\,(2b\,\mathrm da-a\,\mathrm db)}{r\,(a^2+b)},\qquad a=xz\ (l_7),\quad a=xy\ (l_8).
\end{equation}
Writing $\vec J=\sum_{w\ge0}\eps^w\vec J^{(w)}$, the weight-0 part $\vec J^{(0)}$ is constant (\texttt{sol0.m})
and $\mathrm d\vec J^{(w)}=A\,\vec J^{(w-1)}$.

\section{The method of Czakon and Tancredi}
We follow ref.~\cite{CT} exactly (sections 2 and 4 there).

\paragraph{One system, square roots included.}
The unknowns are the $\eps$-coefficients of weight $1,\dots,W$ and the square root, stacked into one vector
\begin{equation}
 F=\big(\vec J^{(1)},\dots,\vec J^{(W)},\,r\big)\in\mathbb C^{371W+1},
\end{equation}
which obeys the first-order (non-linear, through $r$) system
\begin{equation}\label{eq:system}
 \frac{\mathrm d\vec J^{(w)}}{\mathrm dt}=A_t(r)\,\vec J^{(w-1)}\quad(w=1,\dots,W),\qquad
 \frac{\mathrm dr}{\mathrm dt}=\frac{r}{2b}\frac{\mathrm db}{\mathrm dt},
\end{equation}
where $A_t=\sum_k a_k\,\frac{\mathrm d}{\mathrm dt}\log l_k$ is the connection contracted with the tangent of the path,
and eq.~\eqref{eq:odd} is evaluated with the \emph{current value of the component $r$}.  The square root is therefore never
evaluated: its branch is fixed once, at the boundary point, and is then continued by the differential equation
itself (eq.~(2.6) of \cite{CT}).  This is what makes the analytic continuation automatic.
For $W=6$ the system has $2227$ components (for $W=2$: $743$).

\paragraph{Path.}
From the boundary point $\vec x_0$ to the target $\vec x_1$, along the complexified straight line (eq.~(4.1) of \cite{CT})
\begin{equation}
 z_k(t)=x_{0k}+\big(1+4i\delta_k(1-t)\big)\big(x_k(t)-x_{0k}\big),\qquad x(t)=x_0+t(x_1-x_0),\qquad t\in[0,1],
\end{equation}
with $\vec\delta=(0.1,0.2)$ (as in \cite{CT}).

\paragraph{Integrator.}
\texttt{boost::numeric::odeint::bulirsch\_stoer}: modified-midpoint steps with $n_k=2(k+1)$ sub-steps,
Richardson extrapolation in $h^2$ up to order $k_{\max}=8$, adaptive step size and extrapolation order.
The local-error estimate is the default of \cite{CT}: the maximal \emph{absolute} error of all components
(\texttt{bulirsch\_stoer(eps\_abs = error, eps\_rel = 0)}), required to be $\le$ \texttt{error}.
Floating-point types: \texttt{double} and \texttt{dd\_real} (double-double, $\simeq 32$ digits, QD library).
Boost~1.83's \texttt{bulirsch\_stoer.hpp} does not compile with \texttt{dd\_real} (ambiguous
\texttt{size\_t}$\to$\texttt{dd\_real} conversions); a local copy only adds explicit \texttt{double} casts of the
integer sequences, the algorithm is unchanged.  Double-double arithmetic must be compiled without FMA contraction
(\texttt{-ffp-contract=off}), otherwise the error-free transformations of QD break and the extrapolation never converges.

\paragraph{Code.}
\texttt{diffeqs.hpp} reproduces the interface of \cite{CT} (\texttt{diffeqs<Real>}, \texttt{connection\_t},
\texttt{vector\_field\_t}, \texttt{path\_t}, boundary stream, \texttt{evaluate}, exceptions for max steps/evaluations/time and
min step size).  \texttt{h\_family.cpp} supplies the connection (14 letters, the $a_k$ assembled on the union sparsity pattern,
43\,410 non-zeros) and the vector field for $r$.  The matrix entries are read with 40 significant digits.

\section{Boundary values}
\subsection{Euclidean region: corner 1}
The available boundary information are the shuffle-regularised constants $\vec c_w$ at the corner $x,y\to0$
(coefficient of $\log^0x\log^0y$).  Near the corner
\begin{equation}
 \vec J^{(w)}=\sum_{a+b\le w}\log^a x\,\log^b y\;\vec C^{(w)}_{ab}(x,y),\qquad
 \vec C^{(w)}_{ab}(0,0)=\frac{a_1^{\,a}\,a_2^{\,b}}{a!\,b!}\,\vec c_{w-a-b} ,
\end{equation}
($a_1,a_2$: matrices of the letters $x$, $y$; the other letters vanishing at the corner, $x+y$ and $x-y+\dots$,
annihilate all $\vec C_{ab}(0,0)$).  On the ray $x=sX$, $y=sY$ the DE becomes
$\mathrm d\vec J/\mathrm ds=\eps\,(R/s+\sum_{i\ge0}B_i s^i)\vec J$ with $R=\sum_k m_k a_k$ ($m_k$ the order of vanishing of
$l_k$) and $B_i=\sum_k B_{k,i}\,a_k$, $B_{k,i}$ the Taylor coefficients of the regular part of $\partial_s\log l_k$.
The solution is a log-power series
\begin{equation}
 \vec J^{(w)}(s)=\sum_{k=0}^{w}\log^k s\sum_{i=0}^{N}\vec P^{(w)}_{k,i}\,s^i ,
\end{equation}
obtained weight by weight: with $\vec G_{k,i}=\sum_q a_q\big(m_q\vec P^{(w-1)}_{k,i}+\sum_{j<i}B_{q,j}\vec P^{(w-1)}_{k,i-1-j}\big)$
(the coefficient of $s^{i-1}\log^k s$ in $A_s\vec J^{(w-1)}$),
\begin{equation}
 \int\! s^{-1}\log^k s=\frac{\log^{k+1}s}{k+1},\qquad
 \int\! s^{i-1}\log^k s=s^i\sum_{j=0}^{k}\frac{(-1)^j k!}{(k-j)!\,i^{j+1}}\log^{k-j}s\quad(i\ge1),
\end{equation}
and the integration constant fixed by the $s$-regularised ray constants
\begin{equation}
 \vec P^{(w)}_{0,0}\ \ni\ \vec c^{\,\rm ray}_w=\sum_{a+b\le w}\vec C^{(w)}_{ab}(0,0)\,\log^aX\,\log^bY .
\end{equation}
The ray constants are computed exactly (rational matrices times the constants $\pi,\zeta_3,\zeta_5,\dots$) and rounded
to 36 digits; the $B_{k,i}$ are computed by exact power-series algebra at working precision
($q'/q$ for the plain letters, eq.~\eqref{eq:odd} with $r=s\sqrt{XY}\sqrt{1-s(X+Y)}$ for the odd ones).
We use $(X,Y)=(1,0.7)$, $s_0=0.1$, $N=60$ (truncation error far below $10^{-32}$), i.e.\ $\vec x_0=(0.1,0.07)$, and as an
independent second start $(X,Y)=(0.4,1)$, $s_0=0.12$.  The boundary is produced in double-double
(\texttt{frobenius.cpp}).  Against the exact weight-1,2 solution the boundary vectors agree to
$\simeq10^{-26}$ ($w=1$) and $\simeq10^{-25}$ ($w=2$).
(The earlier double-precision version of this step, with FFT-computed $B_{k,i}$, was only good to $10^{-9}$.)

\subsection{Continued (scattering) region}
The region $x,y<0$, $z>1$ is parametrised as $x=-w/v$, $y=-u/v$, $z=1/v$, $u+w+v=1$.  Its corner
$u,w\to0$ (CC1) is the same corner approached from the negative quadrant.  The boundary constants
$\vec c^{\,\rm CC1}_w$ (\texttt{transport.boundaries/H/dist/bdy\_CC1\_w*.m}) are the corner-1 data continued with
$\log x\to\log w-\log(1-w-u)+i\pi$, $\log y\to\log u-\log(1-w-u)+i\pi$, and are regularised in $\log w,\log u$.
Since $\log v=O(u,w)$, the same construction applies with $\vec c\to\vec c^{\,\rm CC1}$, a ray with $X,Y<0$, and
$\log X,\log Y\to\log|X|,\log|Y|$.  The continuation $x\to x+i0$, $y\to y+i0$ (both phases $+\pi$) rotates $b=xyz$ by
$2\pi$, hence
\begin{equation}
 r=-\sqrt{|xyz|}\qquad\text{in the continued region,}
\end{equation}
which is the only change in the series of the odd letters and in the boundary value of $r$.
Starts: $(X,Y)=(-1,-0.7)$, $s_0=0.1$ and $(X,Y)=(-0.4,-1)$, $s_0=0.12$.

\section{Validation}
\paragraph{Exact weight $\le2$.}
The exact solution through weight 2 is a combination of GPLs with letters $\{0,1\}$ in $y$ and $\{0,-y,1-y\}$ in $x$
(no square roots).  In the continued region it is evaluated on the real negative axes with principal branches,
$G(0;x)=\log|x|+i\pi$, i.e.\ $x+i0$, $y+i0$; the weight-2 GPLs with non-zero indices are computed by direct
numerical integration along $[0,x]$, on which no index lies (this avoids spurious cuts of closed $\mathrm{Li}_2$ formulas).
Digits are $-\log_{10}|J_{\rm num}-J_{\rm ex}|/|J_{\rm ex}|$ over the components with $|J_{\rm ex}|\ge10^{-4}$.

\paragraph{Weights 3--6.}
As recommended in \cite{CT}: integration from two different boundary points, and, for the continued region, from
the Euclidean boundary along a path crossing $x=0$, $y=0$ above the origin (so that the system itself continues
$r$ and the logarithms), which tests the transported CC1 constants independently.

%RESULTS%

\begin{thebibliography}{9}
\bibitem{CT} M.~Czakon and L.~Tancredi, \emph{Solution of Canonical Differential Equations for Integrals on Arbitrary
Geometries}, arXiv:2606.30354.
\end{thebibliography}
\end{document}
